\section{Explicit certificates for the positive subclasses}
\label{app:positive-certificates}

This appendix supplies the reduction and the coefficient formulas used
in Proposition~\ref{prop:same-symmetry}. We first compute the
commutator, then rearrange its window terms by Pl\"ucker identities.
All index sets are finite; a sum over an empty set is zero.

\subsection{Reduction for pairs with the same symmetry}
\label{app:same-symmetry-reduction}

For symmetric Toeplitz matrices write $X_{ij}=x_{|i-j|}$ and
$Y_{ij}=y_{|i-j|}$. For skew-symmetric Toeplitz matrices write
$X_{ij}=\operatorname{sgn}(i-j)x_{|i-j|}$ and similarly for $Y$,
with $x_0=y_0=0$. In both cases $z_{ab}=x_ay_b-x_by_a$.
In the symmetric case the diagonal multiplicities are $w_0=N$ and
$w_a=2(N-a)$ for $a>0$. The weighted Lagrange identity
\cite[Lemma~10]{laszlo2012} gives
\[
 2\norm X_F^2\norm Y_F^2-2\langle X,Y\rangle_F^2
 =4N\sum_{b=1}^{N-1}(N-b)z_{0b}^2
   +8\sum_{1\le a<b<N}(N-a)(N-b)z_{ab}^2.
\]
The same formula holds in the skew-symmetric case with its first sum
equal to zero.

Use matrix indices $0,\ldots,N-1$, set $m=N-1$, and fix
$\delta=j-i>0$. For a symmetric pair, splitting the summation index
in $[X,Y]_{ij}$ at $i$ and $j$ gives
\[
 [X,Y]_{i,i+\delta}
 =\sum_{t=0}^{i}z_{t,t+\delta}
  -\sum_{t=0}^{m-\delta-i}z_{t,t+\delta}.
\]
The terms with an index strictly between $i$ and $j$ cancel in pairs.
For a skew-symmetric pair the displayed expression changes sign.
Both commutators are skew-symmetric. Each nonempty window occurs
twice above the diagonal, so
\[
 \norm{[X,Y]}_F^2
 =4\sum_{\delta=1}^{m-1}
   \sum_{\substack{k\ge0\\2k<m-\delta}}
   \left(\sum_{t=k+1}^{m-\delta-k}z_{t,t+\delta}\right)^2.
\]
Reflection $t=N-\delta-p$ turns each window into the corresponding
sum of $\hat z_{p,p+\delta}=z_{N-p-\delta,N-p}$. Combining these
identities proves \eqref{eq:positive-reduction}.

\subsection{The window certificate}\label{app:positive-windows}

For $1\le\delta<m$, put $h_\delta=m-\delta$ and define
\[
 \kappa_\delta(p,q)
 =\min(p,q,h_\delta+1-p,h_\delta+1-q),
 \qquad 1\le p,q\le h_\delta.
\]
This counts the nested intervals
$[k+1,h_\delta-k]$ containing both $p$ and $q$. Consequently,
\begin{equation}\label{eq:positive-tent}
 \sum_{\substack{k\ge0\\2k<h_\delta}}
 \left(\sum_{p=k+1}^{h_\delta-k}v_p\right)^2
 =\sum_{p,q=1}^{h_\delta}\kappa_\delta(p,q)v_pv_q.
\end{equation}
Thus the uncorrected Gram of $\Psi_m$ has diagonal
$2p(p+\delta)-\kappa_\delta(p,p)$ and entries
$-\kappa_\delta(p,q)$ between distinct wedges of gap $\delta$.

For $3\le\sigma\le2m-1$, set
\[
 \begin{gathered}
 I_\sigma=\{a\in\mathbb Z:\max(1,\sigma-m)\le a<\sigma/2\},\\
 c_\sigma=\max(0,\sigma-m-1),\qquad
 L_\sigma=\min(\sigma,m+1),\\
 \mathcal W_k^\sigma
 =\{a\in I_\sigma:c_\sigma+1+k\le a\le L_\sigma-1-k\},\\
 t_\sigma(a,b)
 =\min(a-c_\sigma,b-c_\sigma,L_\sigma-a,L_\sigma-b).
 \end{gathered}
\]
For $1\le p<q\le h_\delta$ with $e=q-p\le\delta$, define
\[
 \omega_\delta(p,q)=
 \begin{cases}
 \kappa_\delta(p,q)+\kappa_e(p,p+\delta),&e<\delta,\\
 \kappa_\delta(p,q),&e=\delta.
 \end{cases}
\]
The remaining diagonal coefficient is
\begin{equation}\label{eq:positive-slack}
 s_\delta(p)=2p(p+\delta)-\kappa_\delta(p,p)
 -t_{2p+\delta}(p,p)
 -\sum_{\substack{1\le q\le h_\delta\\0<|q-p|\le\delta}}
 \omega_\delta(\min(p,q),\max(p,q)).
\end{equation}

\begin{proposition}[An explicit window decomposition]
\label{prop:positive-window-certificate}
For every $m\ge1$, the reflected wedges in
\eqref{eq:positive-reduced-form} satisfy
\begin{equation}\label{eq:positive-window-certificate}
 \begin{aligned}
 \Psi_m(\hat z)
 &=\sum_{\sigma=3}^{2m-1}\sum_{k\ge0}
     \left(\sum_{a\in\mathcal W_k^\sigma}
                  \hat z_{a,\sigma-a}\right)^2 +\sum_{\delta=1}^{m-1}
   \sum_{\substack{1\le p<q\le h_\delta\\q-p\le\delta}}
   \omega_\delta(p,q)
       (\hat z_{p,p+\delta}-\hat z_{q,q+\delta})^2 +\sum_{\delta=1}^{m-1}\sum_{p=1}^{h_\delta}
      s_\delta(p)\hat z_{p,p+\delta}^2.
 \end{aligned}
\end{equation}
All nonempty edge weights are positive integers, and
$s_\delta(p)\ge p^2$. In particular, the right-hand side is SoS
in the original symbol variables.
\end{proposition}

\begin{proof}

For $p\ge1$, $1\le\alpha<\beta$ and
$p+\alpha+\beta\le m$, the Pl\"ucker identity is
\[
 R_{p,\alpha,\beta}
 :=\hat z_{p,p+\alpha}\hat z_{p+\beta,p+\alpha+\beta}
  -\hat z_{p,p+\beta}\hat z_{p+\alpha,p+\alpha+\beta}
  +\hat z_{p,p+\alpha+\beta}\hat z_{p+\alpha,p+\beta}=0.
\]
Add $2\kappa_\alpha(p,p+\beta)R_{p,\alpha,\beta}$ to
$\Psi_m$ for every such triple. These additions preserve the
polynomial. Their effect on the Gram entries is as follows.
For two wedges of the same gap $\delta$ at positions $p<q$,
write $e=q-p$. If $e>\delta$, the correction cancels the entry
$-\kappa_\delta(p,q)$. If $e<\delta$, it changes that entry to
$-\kappa_\delta(p,q)-\kappa_e(p,p+\delta)=-\omega_\delta(p,q)$.
If $e=\delta$, no four distinct indices are present and the entry
remains $-\omega_\delta(p,q)$.

The third product in $R_{p,\alpha,\beta}$ couples wedges on the
same antidiagonal. For $a<b$ in $I_\sigma$, its coefficient in
the Gram is
\[
 \kappa_{b-a}(a,\sigma-b)
 =\min(a,\sigma-b,m+1-b,m+1+a-\sigma)
 =t_\sigma(a,b).
\]
Counting the sets $\mathcal W_k^\sigma$ containing both $a$ and
$b$ expresses these entries by the first sum in
\eqref{eq:positive-window-certificate}. The second sum gives
the negative same-gap entries, and \eqref{eq:positive-slack}
makes the diagonals agree. All remaining entries are zero.
This proves the identity.

It remains to check the signs. Each tent value in its range is
a positive integer, so $\omega_\delta(p,q)\ge1$.
The bounds $\kappa_\delta(p,p)\le p$ and
$t_{2p+\delta}(p,p)\le p$ are immediate.
For a right neighbour at distance $e<\delta$ the edge weight
is at most $2p$, and at distance $e=\delta$ it is at most $p$.
For a left neighbour these bounds are $2(p-e)$ and $p-\delta$,
respectively. Including missing boundary neighbours only increases
the upper bound, whence
\[
 \sum_q\omega_\delta(\min(p,q),\max(p,q))
 \le
 \begin{cases}
 (2\delta-1)p+p(p-1),&p\le\delta,\\
 (4\delta-2)p-\delta^2,&p>\delta.
 \end{cases}
\]
Inserting these estimates into \eqref{eq:positive-slack} gives
\[
 s_\delta(p)\ge
 \begin{cases}
 p^2,&p\le\delta,\\
 p^2+(p-\delta)^2,&p>\delta.
 \end{cases}
\]
The sums in the certificate are finite because the windows are
eventually empty. Every wedge is quadratic in the original variables,
so the displayed identity is a finite sum of homogeneous quadratic
squares. Integer weights may equivalently be realized by repeating
the corresponding squares.
\end{proof}

\subsection{The coupling when both factors have both parts}
\label{app:positive-coupling}

Let $X=S_1+K_1$ and $Y=S_2+K_2$, with $S_1,S_2$ symmetric and $K_1,K_2$
skew-symmetric. Expanding the two Frobenius-orthogonal parts of
the commutator gives
\[
 F(X,Y)=F(S_1,S_2)+F(K_1,K_2)+F(S_1,K_2)+F(K_1,S_2)+\mathcal C,
\]
where
\[
 \begin{aligned}
 \mathcal C= -4\langle S_1,S_2\rangle_F\langle K_1,K_2\rangle_F
 -2\langle[S_1,S_2],[K_1,K_2]\rangle_F -2\langle[S_1,K_2],[K_1,S_2]\rangle_F.
 \end{aligned}
\]
For Toeplitz matrices the four displayed $F$ terms have the
certificates of Section~\ref{sec:positive-subclasses}. The term
$\mathcal C$ vanishes if one factor has only one symmetry part,
but is present in general. No sign or SoS assertion for
$\mathcal C$ is needed in the one-sided proof.
